% =====================================================================
%  Pregled teme master rada — kratak dokument za prvi uvid mentora.
%  Build: lualatex pregled-teme.tex (dva puta, zbog brojeva tabela)
%  Brojevi u tabelama 1–3 odgovaraju stanju repozitorijuma na dan pisanja.
% =====================================================================
\documentclass[11pt,a4paper]{article}

\usepackage[margin=2.5cm]{geometry}
\usepackage{fontspec}
\setmainfont{Times New Roman}
\setmonofont{Consolas}[Scale=0.88]
\usepackage[main=serbian]{babel}
\usepackage{setspace}
\setstretch{1.25}
\setlength{\parindent}{0pt}
\setlength{\parskip}{6pt}

\usepackage{graphicx}
\graphicspath{{../thesis/slike/}}
\usepackage{booktabs}
\usepackage{tabularx}
\usepackage{array}
\newcolumntype{Y}{>{\raggedright\arraybackslash}X}
\newcolumntype{L}[1]{>{\raggedright\arraybackslash}p{#1}}
\usepackage{enumitem}
\setlist{itemsep=3pt}
\usepackage{xcolor}
\usepackage{caption}
\captionsetup{labelsep=colon}
\usepackage{titlesec}
\titleformat{\section}{\normalfont\Large\bfseries}{\thesection}{1em}{}
\titleformat*{\subsection}{\normalfont\large\bfseries}
\setcounter{secnumdepth}{1}  % podnaslovi u odjeljku 4 bez broja, kao u primjeru
\usepackage{fancyhdr}
\pagestyle{fancy}
\fancyhf{}
\fancyfoot[R]{\thepage}
\renewcommand{\headrulewidth}{0pt}
\fancypagestyle{plain}{\fancyhf{}\fancyfoot[R]{\thepage}\renewcommand{\headrulewidth}{0pt}}
\usepackage[hidelinks]{hyperref}

\addto\captionsserbian{\renewcommand{\tablename}{Tabela}}
\newcommand{\kod}[1]{\texttt{#1}}

\begin{document}

% --- Naslovni blok -------------------------------------------------------
\begin{center}
  \includegraphics[height=1.6cm]{logo-ucg.png}\par\vspace{1.2em}
  UNIVERZITET CRNE GORE\\
  ELEKTROTEHNIČKI FAKULTET\par\vspace{1.6em}
  {\LARGE\bfseries Pregled teme master rada\par}\vspace{1.2em}
  {\large Dizajn i implementacija jednostavnog asemblerskog prevodioca\\
          za prostu arhitekturu (podskup RISC-V RV32I)\par}\vspace{1.4em}
  Milica Delibašić\par\vspace{0.8em}
  Podgorica, oktobar 2026.
\end{center}

\vspace{0.5em}
\noindent\rule{\textwidth}{0.6pt}
{\itshape Ovaj dokument je kratak pregled teme, pripremljen radi prvog uvida. Nije skraćena
verzija rada: sadrži ideju, ono što je do sada implementirano i provjereno, predloženu
strukturu rada i spisak onoga što ostaje.\par}
\noindent\rule{\textwidth}{0.6pt}

% =====================================================================
\section{Tema i problem}

Tema je projektovanje i implementacija jednostavnog \emph{dvoprolaznog asemblera}: programa
koji prihvata izvorni kod napisan na asemblerskom jeziku i prevodi ga u mašinski kod. Za ciljnu
arhitekturu izabran je osnovni cjelobrojni skup instrukcija RISC-V arhitekture (RV32I), otvorene
arhitekture sa malim brojem pravilnih formata instrukcija, koja se sve više koristi i u industriji
i u nastavi.

Polazno zapažanje je sljedeće. Industrijski asembleri, poput GNU asemblera, pouzdano prevode
programe, ali su veliki i složeni alati čija unutrašnja struktura nije pogodna za upoznavanje sa
principima prevođenja. Obrazovni simulatori su usmjereni na \emph{izvršavanje} programa, a sam
postupak \emph{prevođenja} ostaje u pozadini:

\begin{itemize}
  \item nije vidljivo kako asembler rješava \emph{reference unaprijed}, tj.\ upotrebu labele
        koja je definisana tek kasnije u programu;
  \item kodiranje instrukcije u bitska polja, posebno kod grana i skokova, gdje su bitovi
        pomjeraja ispremještani, prikazuje se samo kao gotova heksadecimalna riječ;
  \item faze prevodioca iz teorije (leksička, sintaksna, semantička analiza i generisanje koda)
        nisu jasno razdvojene u kodu postojećih alata.
\end{itemize}

Problem koji rad obrađuje je, dakle, kako napraviti asembler koji je dovoljno mali da se može
u potpunosti opisati i razumjeti, a dovoljno potpun da prevodi smislene programe (labele, podaci,
grananja, petlje, pristup memoriji), uz jasne poruke o greškama i uz mašinski kod koji se može
provjeriti poređenjem sa referentnim asemblerom.

% =====================================================================
\section{Istraživačka pitanja}

\textbf{Glavno pitanje.} Na koji način se može projektovati i implementirati jednostavan
dvoprolazni asemblerski prevodilac za podskup RISC-V arhitekture (RV32I) tako da generiše
ispravan mašinski kod i daje razumljive dijagnostičke poruke o greškama?

\textbf{P1.} Koji minimalni skup instrukcija, direktiva i pseudo-instrukcija RV32I arhitekture
je dovoljan za pisanje smislenih programa (aritmetika, grananje, petlje, pristup memoriji)?

\textbf{P2.} Kako se klasične faze prevodioca preslikavaju na arhitekturu dvoprolaznog asemblera
i koje su posebnosti obrade referenci unaprijed kod grana i skokova?

\textbf{P3.} U kojoj mjeri se mašinski kod razvijenog asemblera poklapa sa kodom referentnog
asemblera (GNU \kod{as}) za isti ulazni program?

% =====================================================================
\section{Šta je do sada implementirano}

Rješenje se razvija u jeziku Python~3.12, bez spoljnih biblioteka, kao alat komandne linije
(\kod{python -m asm program.asm -o out/}). Svaka faza prevodioca je zaseban modul.
Implementirani su leksički analizator, tabela instrukcija i registara, kodiranje svih šest
formata instrukcija sa provjerom opsega, tabela simbola, zapis izlaznih fajlova i komandna
linija. Parser i dva prolaza su projektovani (algoritam i gramatika su definisani), ali još nisu
napisani.

\begin{table}[h]
  \centering
  \caption{Obim podržanog instrukcijskog skupa}
  \begin{tabularx}{0.82\textwidth}{@{}lXr@{}}
    \toprule
    \textbf{Format} & \textbf{Grupe instrukcija} & \textbf{Broj} \\
    \midrule
    R & aritmetika, logika, pomjeranja i poređenja nad registrima & 10 \\
    I & aritmetika i logika sa konstantom, pomjeranja za konstantu,
        učitavanje iz memorije, skok kroz registar & 15 \\
    S & upis u memoriju & 3 \\
    B & uslovna grananja & 6 \\
    U & gornja konstanta (\kod{lui}, \kod{auipc}) & 2 \\
    J & bezuslovni skok (\kod{jal}) & 1 \\
    \midrule
    \multicolumn{2}{@{}l}{Nativnih instrukcija ukupno (implementirano)} & 37 \\
    \multicolumn{2}{@{}l}{Pseudo-instrukcija (\kod{nop mv not j ret li la}; planirano)} & 7 \\
    \multicolumn{2}{@{}l}{Direktiva asemblera (\kod{.text .data .word .byte} \dots; planirano)} & 7 \\
    \bottomrule
  \end{tabularx}
\end{table}

\begin{table}[h]
  \centering
  \caption{Stanje po modulima}
  \begin{tabularx}{0.9\textwidth}{@{}llX@{}}
    \toprule
    \textbf{Modul} & \textbf{Stanje} & \textbf{Uloga} \\
    \midrule
    \kod{lexer.py}    & gotovo   & tekst $\rightarrow$ tokeni, sa linijom i kolonom svakog tokena \\
    \kod{isa.py}      & gotovo   & tabela instrukcija (format, opcode, funct3, funct7) i registara \\
    \kod{encode.py}   & gotovo   & pakovanje polja u 32-bitnu riječ za formate R/I/S/B/U/J \\
    \kod{symbols.py}  & gotovo   & tabela simbola, dupla i nedefinisana labela \\
    \kod{output.py}   & gotovo   & binarni (little-endian) i heksadecimalni izlaz \\
    \kod{cli.py}      & gotovo   & komandna linija \\
    \kod{parser.py}   & sljedeće & tokeni $\rightarrow$ labele, direktive, instrukcije \\
    \kod{assembler.py}& planirano& prvi prolaz (adrese, simboli) i drugi prolaz (kodiranje) \\
    \bottomrule
  \end{tabularx}
\end{table}

% =====================================================================
\section{Ideja na konkretnom primjeru}

Pošto je cilj rada da postupak prevođenja bude jasan, ideja se najkraće objašnjava primjerom.
Riječi mašinskog koda u ovom odjeljku proizvodi već implementirani modul za kodiranje.

\subsection{Dva prolaza i reference unaprijed}

\begin{center}
\begin{minipage}[t]{0.3\textwidth}
\small
\begin{verbatim}
    addi t0, x0, 0
    addi t1, x0, 10
loop:
    addi t0, t0, 1
    blt  t0, t1, loop
\end{verbatim}
\end{minipage}%
\hfill
\begin{minipage}[t]{0.66\textwidth}
\small
\begin{tabular}{@{}lll@{}}
  \toprule
  \textbf{Adresa} & \textbf{Instrukcija} & \textbf{Riječ} \\
  \midrule
  \kod{0x00} & \kod{addi t0, x0, 0}   & \kod{0x00000293} \\
  \kod{0x04} & \kod{addi t1, x0, 10}  & \kod{0x00A00313} \\
  \kod{0x08} & \kod{addi t0, t0, 1}   & \kod{0x00128293} \\
  \kod{0x0C} & \kod{blt t0, t1, loop} & \kod{0xFE62CEE3} \\
  \bottomrule
\end{tabular}
\end{minipage}
\end{center}

U \emph{prvom prolazu} asembler samo računa adrese: brojač lokacija raste za 4 po instrukciji, a
labela \kod{loop} upisuje se u tabelu simbola sa adresom \kod{0x08}. Ništa se još ne prevodi.
U \emph{drugom prolazu}, kada su sve adrese poznate, svaka instrukcija se kodira. Za granu
\kod{blt} upisuje se \emph{razlika} adresa, $\kod{0x08} - \kod{0x0C} = -4$, a ne sama adresa
labele (PC-relativno adresiranje). Da je labela definisana tek poslije grane, postupak bi bio isti,
jer je prvi prolaz već završen.

\subsection{Rastavljanje instrukcije na bitska polja}

Instrukcija \kod{add x3, x1, x2} kodira se u riječ \kod{0x002081B3}, po R formatu:

\begin{center}
\small
\begin{tabular}{@{}cccccc@{}}
  \toprule
  \kod{funct7} & \kod{rs2} & \kod{rs1} & \kod{funct3} & \kod{rd} & \kod{opcode} \\
  \midrule
  \kod{0000000} & \kod{00010} & \kod{00001} & \kod{000} & \kod{00011} & \kod{0110011} \\
  0 & 2 (\kod{x2}) & 1 (\kod{x1}) & 0 & 3 (\kod{x3}) & \kod{0x33} \\
  \bottomrule
\end{tabular}
\end{center}

Kod grane \kod{blt t0, t1, loop} pomjeraj $-4$ zapisuje se kao 13-bitni broj, a njegovi bitovi se
raspoređuju po B formatu na neuzastopna mjesta u riječi:

\begin{center}
\small
\begin{tabular}{@{}cccccccc@{}}
  \toprule
  \kod{imm[12]} & \kod{imm[10:5]} & \kod{rs2} & \kod{rs1} & \kod{funct3} & \kod{imm[4:1]} & \kod{imm[11]} & \kod{opcode} \\
  \midrule
  \kod{1} & \kod{111111} & \kod{00110} & \kod{00101} & \kod{100} & \kod{1110} & \kod{1} & \kod{1100011} \\
  \bottomrule
\end{tabular}
\end{center}

\subsection{Dijagnostika sa tačnom lokacijom}

Svaka greška nosi fajl, liniju i kolonu i prijavljuje se u fazi u kojoj nastaje. Primjeri stvarnog
izlaza trenutne implementacije:

\begin{center}
\begin{minipage}{0.92\textwidth}
\small
\begin{verbatim}
primjer.asm:1:14: lexical error: invalid integer literal '12abc'
primjer.asm:1:9: lexical error: unterminated string literal
primjer.asm:4:5: semantic error: immediate 5000 out of range (-2048..2047)
primjer.asm:9:5: semantic error: branch offset 4100 out of range (-4096..4095)
primjer.asm:12:5: semantic error: shift amount 40 out of range (0..31)
\end{verbatim}
\end{minipage}
\end{center}

% =====================================================================
\section{Verifikacija}

Ispravnost se provjerava automatizovanim testovima (\kod{pytest}). Očekivane vrijednosti
kodiranja izvedene su nezavisno, iz definicije bitskih polja u specifikaciji RISC-V, a ne iz
izlaza same implementacije, da bi test mogao otkriti grešku, a ne samo promjenu ponašanja.
Za grane je svaki od 4096 mogućih pomjeraja kodiran i zatim vraćen nezavisnim dekoderom.

\begin{table}[h]
  \centering
  \caption{Automatizovani testovi}
  \begin{tabularx}{0.85\textwidth}{@{}lrY@{}}
    \toprule
    \textbf{Skup} & \textbf{Broj} & \textbf{Predmet provjere} \\
    \midrule
    Leksička analiza            & 50  & tokeni, brojevi, stringovi, lokacije, greške \\
    Tabela instrukcija          & 74  & polja prema specifikaciji, jedinstvenost kodiranja, registri \\
    Kodiranje instrukcija       & 85  & poznate riječi svih šest formata, granice opsega \\
    Paket i komandna linija     & 8   & izlazni fajlovi, poruke, tabela simbola \\
    \midrule
    Ukupno                      & 217 & svi prolaze \\
    \bottomrule
  \end{tabularx}
\end{table}

Kada parser i dva prolaza budu završeni, dodaju se integracioni testovi nad kompletnim
programima (petlja, grananje, rad sa memorijom, poziv funkcije) i poređenje mašinskog koda sa
izlazom GNU asemblera (\kod{riscv64-unknown-elf-as}) za iste programe.

% =====================================================================
\section{Predložena struktura rada i trenutni status}

Struktura rada prati uputstvo za temu, a izgled je usklađen sa odobrenim master radovima
Elektrotehničkog fakulteta (margine, font, prored, numeracija, redoslijed uvodnih strana).
Pisani dio za sada postoji kao skelet: naslovi svih poglavlja i sekcija, sa bilješkama o sadržaju.

\begin{table}[h]
  \centering
  \caption{Status po poglavljima}
  \begin{tabularx}{\textwidth}{@{}L{6.2cm}lY@{}}
    \toprule
    \textbf{Poglavlje} & \textbf{Status} & \textbf{Napomena} \\
    \midrule
    Uvod & skelet & predmet, cilj i istraživačka pitanja preuzimaju se iz prijave \\
    1. Osnove računarskih arhitektura i asemblera & skelet & literatura izabrana \\
    2. Teorijske osnove prevodilaca & skelet & literatura izabrana \\
    3. Specifikacija jednostavne arhitekture & djelimično & tabele formata i opkodova gotove \\
    4. Projektovanje asemblerskog jezika & djelimično & gramatika (EBNF) i leksička pravila definisani \\
    5. Dizajn i implementacija & prati kod & sekcije odgovaraju modulima \\
    6. Praktična realizacija i primjeri & čeka kod & poslije drugog prolaza \\
    7. Evaluacija i diskusija & čeka kod & poslije poređenja sa GNU \kod{as} \\
    Zaključak & — & na kraju \\
    Prilozi & djelimično & izvorni kod i tabela opkodova uključuju se automatski \\
    \bottomrule
  \end{tabularx}
\end{table}

% =====================================================================
\section{Šta ostaje}

\begin{enumerate}
  \item \textbf{Parser} --- prevođenje tokena u labele, direktive i instrukcije po definisanoj gramatici.
  \item \textbf{Prvi i drugi prolaz} --- dodjela adresa, tabela simbola, kodiranje,
        proširenje pseudo-instrukcija (\kod{li}, \kod{la}, \kod{j}, \kod{ret}, \dots) i direktive za podatke.
  \item \textbf{Izlaz za sekciju podataka} i završetak komandne linije (opciono: listing asembliranja).
  \item \textbf{Integracioni testovi i poređenje sa GNU asemblerom} --- odgovor na pitanje P3.
  \item \textbf{Pisanje teksta rada} --- teorijska poglavlja paralelno sa kodom, zatim evaluacija i zaključak.
  \item Vaše korekcije na prijavi teme i na predloženoj strukturi rada.
  \item Dopuna podataka: broj indeksa, tačan naziv studijskog programa i podaci o mentoru.
\end{enumerate}

% =====================================================================
\section{Literatura (izbor)}

\begin{enumerate}[label=\arabic*., leftmargin=1.6em]
  \item Aho A. V., Lam M. S., Sethi R., Ullman J. D., \emph{Compilers: Principles, Techniques,
        and Tools}, 2.~izd., Pearson/Addison-Wesley, 2006.
  \item Beck L. L., \emph{System Software: An Introduction to Systems Programming}, 3.~izd.,
        Addison-Wesley, 1997.
  \item Waterman A., Asanović K. (ur.), \emph{The RISC-V Instruction Set Manual, Volume I:
        Unprivileged ISA}, verzija 20191213, RISC-V Foundation, 2019.
  \item Patterson D. A., Hennessy J. L., \emph{Computer Organization and Design RISC-V Edition:
        The Hardware/Software Interface}, 2.~izd., Morgan Kaufmann, 2020.
  \item Patterson D., Waterman A., \emph{The RISC-V Reader: An Open Architecture Atlas},
        Strawberry Canyon, 2017.
  \item Waterman A., \emph{Design of the RISC-V Instruction Set Architecture}, doktorska
        disertacija, University of California, Berkeley, 2016.
  \item Salomon D., \emph{Assemblers and Loaders}, Ellis Horwood, 1993.
  \item Nystrom R., \emph{Crafting Interpreters}, Genever Benning, 2021.
  \item Petersen M. B., „Ripes: A Visual Computer Architecture Simulator”, \emph{IEEE/ACM
        Workshop on Computer Architecture Education (WCAE)}, 2021.
  \item Free Software Foundation, \emph{Using as: The GNU Assembler}, dokumentacija GNU Binutils.
\end{enumerate}

\end{document}
